\scanpage{126}
Soit $\deleted{\unsure{\mathcal X}=}(X_i)_{i\in I}$ un système projectif de préschémas, essentiellement affine [i.e. $\exists i_0$, $\forall i\geq i_0$, $\varphi_{i_0i}$ est affine]. Alors $X=\varprojlim_i X_i$ existe, $\mathrm{Top}\,X=\varprojlim_i\mathrm{Top}(X_i)$ $\ldots$.

Soit $x\in X$, on dit que $\unsure{X|x}$ est \uline{épais} si

(Ép) $\exists i$ t.q. $\forall j\geq i$, $x_j$ est un p\textsuperscript{t} maximal de $\varphi_{ij}^{-1}(x_i)$, donc si $k\geq j\geq i$, $\deleted{\unsure{\varphi}}x_k$ est maximal de $\varphi_{jk}^{-1}(x_j)$.

\unsure{Sous quelles conditions est équivalent|Sous quelles conditions est épaisseur}:

(Ép)$'$ $\exists i$, \deleted{t.q.} et une partie ouverte non vide $U_i$ de $\bar x_i$, t.q. $\forall j\geq i$, $\varphi_{ij}^{-1}(U_i)\subset\bar x_j$ (Alors la même propriété est vraie pour $\forall i'\geq i$).

Cela \unsure{prouve} que $\varphi_{ij}^{-1}(U_i)$ est irréductible, de p\textsuperscript{t} générique $x_j$, et donc (Ép)$'$ implique (Ép)\unsure{?}.

Considérons aussi

(Ép)$''$ $\bar x$ \inserted{$=\varprojlim\bar x_i$} est une partie \uline{constructible} de $X$.

Ceci équivaut \inserted{(si $X_i$ sont \unsure{noeths})} à: $\bar x=\varphi_i^{-1}(\bar x_i)$ pour $i$ \unsure{grandement|grand assez} \deleted{(car: pour $i$ \unsure{constructible} $\ldots$)} (donc pour $i'\geq i$)

(Ép)$'''$ $\exists$ une partie constructible \deleted{\unsure{non fermée}} \inserted{ouverte non vide} de $\bar x$, soit $U$.

i.e. $\exists U_i$
\margin{125}
